% !PRINTED 179
\weberpage{179}
%<<CONT
so ist $\nu$ der Grad der Gruppe $T_r$, und wenn die Function $\Theta$, die
wir in \S.\ 43 zur Lösung des Formenproblems angewandt haben,
durch die Substitutionen von $T_r$ in
\[ \Theta,\ \Theta_1,\ \ldots \Theta_{\nu-1} \]
übergeht, so ist
\[ \Phi_r(t) = (t - \Theta)(t - \Theta_1) \ldots (t - \Theta_{\nu-1}) \]
eine Function von $t$ in dem erweiterten Rationalitätsbereiche
$\Omega_r$; der Körper $\Omega(\Theta)$ ist identisch mit $\Omega_r(\Theta)$. Er ist also ein algebraischer Körper $\mu^{\mathrm{ten}}$ Grades über $\Omega$ und $\nu^{\mathrm{ten}}$ Grades über $\Omega_r$.
Das Formenproblem $\mu^{\mathrm{ten}}$ Grades wird demnach durch Adjunction
eines Radicals auf ein Formenproblem $\nu^{\mathrm{ten}}$ Grades zurückgeführt.

\begin{center}
\S.\ 46.

\gesperrt{Der erweiterte Invariantenbegriff.}
\end{center}

Die relativen Invarianten sind im Grunde ein specieller Fall
eines allgemeineren Begriffes, den wir in ähnlicher Weise, wie
die relativen Invarianten, zur Reduction des Formenproblems
anwenden können.

Wir suchen nach Systemen von homogenen Formen gleichen
Grades der Variablen $x_1, x_2 \ldots x_n$:
\weq{1}{X_1, X_2, \ldots X_m,}
von der Beschaffenheit, dass durch die Anwendungen der Substitutionen der Gruppe $S$ die Functionen $X_1, X_2, \ldots X_m$ ein
System $\Sigma$ von \gesperrt{linearen Substitutionen} erleiden; die Substitutionen $\Sigma$ bilden dann gleichfalls eine Gruppe, und zwar
eine Gruppe, die mit $S$ ein- oder mehrstufig isomorph ist. Ist
$m = 1$, so erhalten wir, wie man sieht, den Begriff der relativen
Invarianten. Wir suchen nun alle Substitutionen von $S$, die
jede einzelne der Functionen (1) ungeändert lassen. Diese Substitutionen bilden eine Gruppe, die wir mit $T$ bezeichnen, und
von der wir nachweisen wollen, dass sie ein Normaltheiler von
$S$ ist. Bezeichnen wir mit $s$ eine Substitution aus $S$, durch die
die $X$ die Substitution $\sigma$ erleiden, so erleiden die $X$ durch $s^{-1}$
die Substitution $\sigma^{-1}$. Ist dann ferner $\tau$ eine Substitution aus $T$,
so bleiben durch $\tau$ die $X$ ungeändert. Demnach erleiden durch
$s^{-1}\tau s$ die $X$ die Substitution $\sigma^{-1}\sigma = 1$, d.\,h. sie bleiben un-
%CONT>>
